\section{Shared storage and bounded lifetimes}
\label{sec:sequence-author-contract}

The distinction between a view and its source is useful in self-joins. A query
can use two aliases of the same row to associate two different searches with
the same sequence. The next chapter shows how aliases distinguish two
searches. At the sequence layer, the two aliases can still
refer to the same backing data.

Consider a row with a long sequence \(S\). One consumer needs \(S[a,b)\)
and another needs \(S[c,d)\). Creating two eager copies could load the complete
sequence twice. Two lazy views instead retain the bounds and identify the
source. A \emph{read session} is the period during which source handles and
the state needed to read them remain valid. Within one read session, BioSeqDB
can share backing state or a scanner when the source identities match exactly.

This is a conditional form of reuse.
\begin{itemize}
  \item Equal text in different rows is not enough to establish shared
        backing. The identity of the stored source must match.
  \item Sharing a scanner does not guarantee that overlapping intervals are
        physically read only once. That depends on the reader and its cache.
  \item Existing eager column caching can also share resident data. Laziness
        additionally permits both consumers to avoid requesting the full
        sequence in the first place.
\end{itemize}

\begin{figure}[H]
\centering
\begin{tikzpicture}[x=1cm,y=1cm,font=\footnotesize,
  view/.style={draw=seqNineBlue,fill=seqNineBlue!10,
    minimum width=3.7cm,minimum height=.75cm,align=center}]
  \node[anchor=west,font=\bfseries] at (0,5.9)
    {(a) Two consumers can refer to one source};
  \draw[draw=black!60,fill=black!4] (3.15,4.77) rectangle (10.2,5.37);
  \draw[seqNineBlue,fill=seqNineBlue!17] (4.2,4.77) rectangle (6.7,5.37);
  \draw[seqNineOrange,fill=seqNineOrange!17] (7.25,4.77) rectangle (9.25,5.37);
  \node[anchor=east] at (2.8,5.07) {source $S$};
  \node at (5.45,5.07) {$[a,b)$};
  \node at (8.25,5.07) {$[c,d)$};
  \node[view] (left) at (3.4,3.8) {first alias: view $S[a,b)$};
  \node[view,draw=seqNineOrange,fill=seqNineOrange!10] (right)
    at (9.8,3.8) {second alias: view $S[c,d)$};
  \draw[seq nine flow,seqNineBlue] (left.north)--(5.45,4.72);
  \draw[seq nine flow,seqNineOrange] (right.north)--(8.25,4.72);
  \node at (6.55,2.94) {same source identity in the same read session};
  \node[anchor=west,font=\bfseries] at (0,2.15)
    {(b) Every read must finish while its owner is alive};
  \draw[seq nine flow] (3.45,-1.25)--(12.15,-1.25) node[right] {time};
  \node[anchor=east] at (3.1,1.25) {read session};
  \draw[fill=black!9,draw=black!60] (3.5,1.0) rectangle (10.8,1.5);
  \node[anchor=east] at (3.1,.35) {borrowed views};
  \draw[seqNineBlue,fill=seqNineBlue!15] (4.3,.1) rectangle (7.3,.6);
  \draw[seqNineOrange,fill=seqNineOrange!15] (5.4,-.68) rectangle (9.45,-.18);
  \node at (5.8,.35) {first consumer};
  \node at (7.425,-.43) {second consumer};
  \draw[dashed,black!60] (10.8,-1.45)--(10.8,1.7);
  \node[anchor=west,align=left] at (11.05,.8)
    {release source\\after both reads};
\end{tikzpicture}
\caption{Aliases do not require duplicate sequence payloads. Sharing requires
matching source identities and compatible ownership. The lower panel shows
two read lifetimes fully inside the lifetime of their read session.}
\label{fig:lazy-lifecycle}
\end{figure}

The non-owning lazy view cannot keep the source alive by itself. A call owner
or PostgreSQL read session holds the source handles, retained values, tree
memory, and temporary allocations. A consumer must finish before that owner
is released. Returning an iterator or saving a view that outlives this owner
would violate the rule even if no symbols had yet been read.

PostgreSQL error handling makes this an explicit engineering requirement.
A database error can transfer control without unwinding C++ stack frames.
Ordinary C++ destructors alone are therefore insufficient for releasing
resources across that boundary.\footcite{postgresql17_cpp_extensibility}
The read session registers cleanup actions that also run on the database's
error path. Tree owners retain the sources and older tree memory they depend
on, and release those dependencies through the same ownership scheme.

Asynchronous work needs a longer-lived owner. Before a worker or accelerator
can outlive the current call, it must have stable input ownership. A packed
buffer and its memory permit remain live until the task completes or the GPU
signals completion. A flag saying that a source can be read by workers does
not itself extend the source's lifetime.

These rules let sequence algorithms request a precise interval without also
having to implement PostgreSQL storage and cleanup. They supply ranks in the
required order, keep the source alive, and allocate a complete interval only
when the consumer requires one. The next chapter uses this sequence interface
inside a scan that can produce several alignment results from one stored row.
\flushbottom
